\section{Introduction}\label{in:section}

An exploration of an Ising interface fixes spins along a random boundary.
To use the unexplored region, one needs estimates that remain valid after
this conditioning. Two difficulties arise. The fixed signs can alternate,
and the exposed boundary can contain many narrow excursions. We give
comparison estimates that retain arbitrary common signs, then replace
the boundary visible from a distant target by one of finitely many
regular signed barriers. The replacement controls the whole target
spin law, including observables that vary with the mesh.

\subsection{Finite graphs and pointwise comparison}
For a finite graph $G=(V,E)$, nonnegative couplings $K_e$, and finite real
fields $h_v$, the Ising law is
\[
 \mu(\sigma)=Z^{-1}\exp\!\left(
   \sum_{\{u,v\}\in E}K_{uv}\sigma_u\sigma_v
       +\sum_{v\in V}h_v\sigma_v\right),
 \qquad \sigma\in\{-1,+1\}^{V}.
\]
Prescribing a spin removes that variable and adds its incident
interactions to the fields of its neighbors. By contrast, deleting a
vertex means removing the vertex and all its incident edges.

Let $I,J,F$ be disjoint, with $I,J$ nonempty, and prescribe common spins
on $F$. Form a comparison graph by deleting $F$, contracting $I$ to one
terminal and $J$ to another, and setting all fields to zero. Let $q_0$
be the probability that the two terminals are connected in its
random-cluster model with cluster weight $2$ and edge parameters
$p_e=1-e^{-2K_e}$. Parallel couplings are added on contraction. Write
$\|\nu-\nu'\|_{\mathrm{TV}}=\sup_A|\nu(A)-\nu'(A)|$.

\begin{theorem}[Finite-graph comparison]\label{in:comparison}
For every two configurations $b,b'$ on $J$ and every configuration $a$
on $I$,
\[
 e^{-4\operatorname{arctanh}q_0}
 \leq
 \frac{\mu(\sigma_I=a\mid\sigma_J=b,\sigma_F)}
      {\mu(\sigma_I=a\mid\sigma_J=b',\sigma_F)}
 \leq e^{4\operatorname{arctanh}q_0}.
\]
Either conditional law may be replaced by an arbitrary mixture of
these laws. The bound is symmetric in the two terminal sets.
\end{theorem}

Lemma~\ref{fd:ratio} proves the theorem. For the square lattice at
\[
 \beta_c=\tfrac12\log(1+\sqrt2),\qquad
 B_r=[-r,r]^2\cap\mathbb Z^2,
\]
dual crossings turn it into a uniform buffered comparison. If
$I\subset z+B_R$ and $J\subset(z+B_{(1+c)R})^c$, the density ratio
is bounded above and below by constants depending only on $c>0$.
If $I\subset z+B_w$, $J\subset(z+B_R)^c$, and $R\ge8w$, its
logarithm is at most $C(w/R)^\alpha$ in absolute value, for an
absolute $\alpha>0$. Common pins may occur anywhere off the terminals.
The comparison also passes through a common conditional randomization,
such as internal Edwards--Sokal bond sampling. Precise geometric and
cross-domain forms appear in Corollaries~\ref{fd:buffered-ratio}
and~\ref{fd:finite-cover}.

\subsection{Exploration across a band}
For the exploration results we use the critical nearest-neighbor Ising
model with zero external field and prescribed boundary spins.
Put $Q_r=\{z:\|z\|_\infty<r\}$. Fix $0<u<v$, a larger square
$Q_L$, and a target at positive distance from the band
$\overline Q_v\setminus Q_u$ and the boundary of $Q_L$.
In the inward construction, reveal the outside of $Q_v$ and explore
towards $Q_u$; the target is a compact subset of $Q_u$. In the outward
construction, reveal $Q_u$ and explore towards $Q_v$; the target lies
between $\partial Q_v$ and a more distant fixed wall. All the
clearances are fixed before the lattice mesh $\delta$ tends to zero.
They can be arbitrarily small positive numbers.
When the two contour levels vary, each ranges over a prescribed open
interval. Their closures are disjoint and preserve all these clearances;
these are the separated level windows.

Choose a fixed diagonal in every lattice face and linearly interpolate
the spins on the resulting triangles. The zero set gives noncrossing
polygonal strands. Explore each strand meeting the starting contour
until its first return to that contour or its first hit of the other
contour, revealing the traversed triangles' spins. Include the starting
contour vertices in the initial reveal. The complete transcript is the
record of this adaptive reveal. The Gibbs property holds conditional on
each transcript, as proved in Lemma~\ref{if:stopping}.

A \emph{regular signed barrier} is a deterministic Jordan contour with
finitely many prescribed plus and minus arcs, together with a lattice
boundary layer following those arcs. Its fresh side is the side containing
the target. A finite descriptor is a function of the band spins with values
in a fixed finite set; one value is allowed to be exceptional.

\begin{theorem}[Finite conditional resolution]\label{in:resolution}
Fix the preceding band geometry and use the critical plus-boundary law
in a larger reference square. For every $\epsilon,\eta>0$, one can
choose contour levels in any prescribed separated subwindows,
a finite descriptor $Z^\delta\in\{0,1,\ldots,m\}$, and regular
signed barriers $\Gamma_1,\ldots,\Gamma_m$ with the following properties
for all sufficiently fine meshes.

The exceptional state $0$ has probability at most $\epsilon$. For every
transcript in a good state $i\in\{1,\ldots,m\}$, insert the plus pins of $\Gamma_i$ and
then its minus pins. The sum of the two resulting total-variation
changes on the target is at most $\eta$. The final target-side law
$Q_{i,\delta}$ is the same for every transcript in that state.
Consequently this conclusion holds simultaneously for all bounded target
observables, including mesh-dependent observables. For any fixed finite
list of bands the descriptors have convergent joint laws.

The finite list of barriers and its positive geometric clearances are
chosen before the mesh threshold. If another law has band density at
most $K$ times the reference density, its exceptional probability is
at most $K\epsilon$.
\end{theorem}

The full statement, including a further inward barrier displacement,
is Theorem~\ref{if:band-transfer}. Density comparison transfers the
probability of the exceptional event. To transfer a conditional target
law, the continuation beyond a common source contour must also be the
same; Corollary~\ref{if:descriptor-averages} states this distinction
precisely.

\subsection{Inputs and proof mechanism}
Critical spatial comparison already follows from FK crossing estimates.
Duminil-Copin, Hongler, and Nolin prove polynomial relative decorrelation
for separated FK events and discuss ratio mixing
\cite[Section~5.3, Proposition~5.11, author-hosted version]{DHN2011}.
Here we formulate comparison directly for full spin configurations on
arbitrary terminal sets of a finite ferromagnetic graph, retaining
arbitrary common pins and fields. The positive-association and cluster
representation framework is due to Fortuin--Kasteleyn--Ginibre, Holley,
and Edwards--Sokal \cite{FKG1971,Holley1974,EdwardsSokal1988};
marginal closure of the lattice condition belongs to the multivariate
total positivity theory of Karlin and Rinott \cite{KarlinRinott1980}.
We give the binary marginalization proof and reduce terminal log-odds
rectangles to the extreme rectangle. Balancing the two terminal fields
then permits the arbitrary-field covariance inequality of Ding, Song,
and Sun \cite[Corollary~1.3]{DSS2023}. That inequality is the imported
two-spin result; the pointwise terminal comparison is the finite-graph
deduction proved here. The FK crossing and one-arm estimates of
\cite{DHN2011} supply its critical geometric consequences.

To control a conditional target law, we construct an ordered coupling
in a doubled spin model. A disagreement that reaches a patch carries
magnetic mass of order $w^{7/8}$ on a nearby counting segment of scale
$w$. Pointwise comparison bounds the effect of changing that patch on
a target at distance $s$. Summing over an admissible cut gives ordered
transmission, including its signed-source and total-variation forms
(Section~\ref{tr:section}). As a separate quantitative consequence of
the spin-correlation limits of Chelkak, Hongler, and Izyurov
\cite{CHI2021}, we prove a mesh-first $L^2$ bound of order $(w/s)^{3/8}$
between normalized magnetizations on nearby contours, and a prediction
bound from a nearby contour (Section~\ref{li:section}). The stopping
construction below uses transmission directly.

Benoist and Hongler prove convergence of critical Ising interfaces to
nested $\mathrm{CLE}_3$ and a simultaneous common-edge approximation
for all macroscopic Ising loops \cite{BH2019}. We use that stronger
approximation to preserve coverage, orientation, and multiplicity under
every noncrossing complete resolution of the lattice interfaces
(Section~\ref{re:section}). At generic stopping levels the part of the
boundary visible from the target is a Jordan base with finitely many
slits (Section~\ref{if:section}).

The main geometric step replaces this slit boundary by regular signed
outlines. Finite graph embeddings and planar extension theorems
\cite{Schaefer2021,Moise1977} provide a topological chart. Narrow ribbons
retain the two boundary sides of each slit and the sign switch at its
tip. The spin-correlation limits of \cite{CHI2021} control magnetic
means in fixed regular upper and lower comparison domains. Small
neighborhoods of roots and tips can be overwritten with common pins at
a vanishing cost on a distant test contour. Away from them, the signs
seen by the common free component order every nearby realized wall
between the regular comparisons. Ordered transmission then turns the
small mean gap into target total variation. For an outward band, a
common bounded weight also permits arbitrary prescribed signs on the
distant wall.

This construction gives one barrier for every realizable transcript
in a neighborhood of the limiting slit datum. Its neighborhood and
mesh threshold are chosen before the target observable. A countable
cover, reduced to finitely many continuity sets carrying arbitrarily
high probability, yields Theorem~\ref{in:resolution}
(Sections~\ref{rg:section}--\ref{de:section}).

The positive width of every comparison collar is a fixed geometric
hypothesis; the constants may depend on it. Pointwise ratios allow
arbitrary common pins, whereas the magnetic mass charts and test
contours used in transmission have the free neighborhoods specified in
Definition~\ref{tr:admissible}. These separate hypotheses will be kept
explicit throughout.
